\subsection{The local term in the canonical formula}\label{sec:P04}
Take $n$ sufficiently large, in particular $n\ge8M$.
We distinguish a consequence of Section~\ref{sec:P02} from one that requires
comparison with the continuous profile (Sections~\ref{sec:P05} and~\ref{sec:P06}).
The local term is
\[
 L_n=-\frac12\sum_{ij}
       [\log(1-4\pi_{ij})-4\log(1-\pi_{ij})].       \tag{LC1}
\]
For $0\le t\le1/8$, the convergent scalar series give
\[
 -\tfrac12[\log(1-4t)-4\log(1-t)]
   =\frac12\sum_{m\ge2}\frac{4^m-4}{m}t^m.
\]
The quadratic term is $3t^2$, and the remaining terms are positive.
Their sum is bounded by
\[
 \frac12\sum_{m\ge3}(4t)^m
       =\frac{32t^3}{1-4t}\le64t^3.
\]
Using $\pi\le M/n$ and $\sum\pi=n$ gives
\[
 0\le L_n-3\sum_{ij}\pi_{ij}^2
       \le64\sum\pi_{ij}^3\le\frac{64M^2}{n}.      \tag{LC2}
\]
Estimate (LC2) is a local expansion on the scale $\pi=O(1/n)$.

\paragraph{From the local term to the continuous profile.}
If $\nu_1$ is the bounded Lipschitz continuous profile and
\[
 \max_{ij}|n\pi_{ij}-\nu_1((i-1/2)/n,(j-1/2)/n)|=O(1/n),
                                                        \tag{LC3}
\]
then
\[
 \sum_{ij}\pi_{ij}^2
       =\int_{[0,1]^2}\nu_1(r,c)^2\,dr\,dc+O(1/n).
\]
Indeed, the difference between the squares of the profiles
is uniformly $O(1/n)$, and the two-dimensional quadrature
of a Lipschitz function has error $O(1/n)$.
By (LC2),
\[
 L_n=3\int_{[0,1]^2}\nu_1^2+O(1/n).               \tag{LC4}
\]
This is only one of the summands in $\log Q_{\rm pred}$. The contributions of the entropy and
the pseudodeterminant are studied in Section~\ref{sec:constante}.

Sections~\ref{sec:P05} and~\ref{sec:P06} prove (LC3) for the canonical maximizer,
so (LC4) applies to it.
